% ======================================================================
\chapter[Уравнение Шредингера для атома водорода]%
        {Уравнение Шредингера для атома водорода}
\label{t:10}
% ======================================================================

\begin{qbox}
\textbf{Билет 10.} Уравнение Шредингера для атома водорода. Волновая функция основного
состояния. Классификация электронных состояний атома водорода (без вывода).
\end{qbox}

\section{Содержание билета}

\textbf{Подтемы:}
\begin{enumerate}
 \item гамильтониан и уравнение Шредингера для электрона в кулоновском поле; приведённая масса;
 \item разделение переменных в сферических координатах; радиальное уравнение, эффективный потенциал;
 \item волновая функция и энергия основного состояния (вывод подстановкой), нормировка,
       радиальное распределение, средние значения;
 \item классификация состояний: квантовые числа $n,l,m$ (и $m_s$), уровни $E_n$, вырождение $n^2$
       \emph{(без вывода)}; обозначения $1s,2p,\dots$
\end{enumerate}

\textbf{Необходимые понятия:} стационарное уравнение Шредингера (\tref{t:06}); оператор
момента импульса и сферические функции (\tref{t:09}); модель Бора (\tref{t:03}); лапласиан
в сферических координатах.

\section{Определения и физические принципы}

\begin{definition}[гамильтониан водородоподобного атома]
В системе центра масс (ядро заряда $Ze$, электрон; приведённая масса $\mu\approx\mel$):
\begin{equation}\label{eq:10:H}
  \op H=-\frac{\hbar^2}{2\mu}\nabla^2-\frac{Ze^2}{r}.
\end{equation}
Учитывается только кулоновское притяжение; спин, релятивистские эффекты и размер ядра не учтены.
\end{definition}

\begin{definition}[квантовые числа атома водорода]
\emph{Главное} $n=1,2,\dots$ (определяет энергию); \emph{орбитальное} $l=0,\dots,n-1$ (квадрат момента
$\hbar^2l(l+1)$); \emph{магнитное} $m=-l,\dots,l$ (проекция $L_z=m\hbar$); \emph{спиновое} $m_s=\pm\frac12$.
\emph{Радиальное} квантовое число $n_r=n-l-1$ — число узлов радиальной функции.
\end{definition}

\emph{Физический смысл.} В центральном поле сохраняются энергия, $L^2$ и $L_z$; их операторы попарно
коммутируют, поэтому стационарные состояния можно нумеровать одновременно значениями $E$, $l$, $m$.
Электрон описывается не орбитой, а распределением плотности вероятности («электронным облаком»).

\emph{Область применимости.} Нерелятивистское приближение: поправки порядка $(Z\alpha)^2$ (тонкая
структура), радиационные поправки (лэмбовский сдвиг), сверхтонкая структура (взаимодействие с
магнитным моментом ядра) — малы и здесь не рассматриваются.

\section{Основная теория}

\subsection{Уравнение Шредингера и разделение переменных}

Стационарное уравнение
\begin{equation}\label{eq:10:sse}
  -\frac{\hbar^2}{2\mu}\nabla^2\psi-\frac{Ze^2}{r}\psi=E\psi .
\end{equation}
Лапласиан в сферических координатах выражается через $\op L^2$ (\tref{t:09}):
\begin{equation}\label{eq:10:lap}
  \nabla^2=\frac{1}{r^2}\pdiff{}{r}\Bigl(r^2\pdiff{}{r}\Bigr)-\frac{\op L^2}{\hbar^2r^2}.
\end{equation}
Поэтому
\(
  \op H=-\frac{\hbar^2}{2\mu}\frac{1}{r^2}\pdiff{}{r}\bigl(r^2\pdiff{}{r}\bigr)+\frac{\op L^2}{2\mu r^2}-\frac{Ze^2}{r}
\)
и решение ищется в виде
\begin{equation}\label{eq:10:sep}
  \psi_{nlm}(r,\theta,\varphi)=R_{nl}(r)\,Y_{lm}(\theta,\varphi).
\end{equation}
Подставив и использовав $\op L^2Y_{lm}=\hbar^2l(l+1)Y_{lm}$, получаем \emph{радиальное уравнение}
\begin{equation}\label{eq:10:radial}
  -\frac{\hbar^2}{2\mu}\frac{1}{r^2}\diff{}{r}\Bigl(r^2\diff{R}{r}\Bigr)
  +\Bigl[\frac{\hbar^2l(l+1)}{2\mu r^2}-\frac{Ze^2}{r}\Bigr]R=ER .
\end{equation}
С подстановкой $R=u/r$ оно сводится к одномерному уравнению
$-\frac{\hbar^2}{2\mu}u''+U_{\text{эфф}}(r)u=Eu$ с граничным условием $u(0)=0$ и
\begin{equation}\label{eq:10:Ueff}
  U_{\text{эфф}}(r)=-\frac{Ze^2}{r}+\frac{\hbar^2l(l+1)}{2\mu r^2}.
\end{equation}
Второй член — \emph{центробежная энергия}: при $l\neq0$ она отталкивает электрон от ядра (рис.~\ref{fig:10:Ueff}).
Связанные состояния существуют при $E<0$.

\begin{figure}[htbp]
\centering
\begin{tikzpicture}
\begin{axis}[width=0.68\linewidth, height=5.4cm, axis lines=middle, xlabel={$r/\aB$}, ylabel={$U_{\text{эфф}}$, Ry},
   xmin=0, xmax=12, ymin=-2.2, ymax=1.3, samples=300, restrict y to domain=-3:2,
   xlabel style={at={(ticklabel* cs:1)}, anchor=west}, ylabel style={at={(ticklabel* cs:1)}, anchor=south},
   legend style={font=\small, draw=none, at={(0.98,0.30)}, anchor=east}]
  \addplot[very thick, accent, domain=0.3:12] {-2/x};
  \addlegendentry{$l=0$}
  \addplot[very thick, keycol, domain=0.3:12] {-2/x + 2/x^2};
  \addlegendentry{$l=1$}
  \addplot[very thick, dercol, domain=0.5:12] {-2/x + 6/x^2};
  \addlegendentry{$l=2$}
  \addplot[gray, dashed, domain=0:12] {-1/4};
  \node[font=\small, anchor=south] at (axis cs:10.5,-0.25) {$E_2$};
\end{axis}
\end{tikzpicture}
\caption{Эффективный потенциал~\eqref{eq:10:Ueff} для $Z=1$ в единицах $\Ry=e^2/2\aB$
($U_{\text{эфф}}/\Ry=-2/\rho+l(l+1)/\rho^2$, $\rho=r/\aB$). Центробежный барьер при $l\neq0$
выталкивает электрон от ядра.}
\label{fig:10:Ueff}
\end{figure}

\subsection{Основное состояние}

Основное состояние сферически симметрично ($l=m=0$) и не имеет узлов. Подстановка $\psi=C\ee^{-r/a}$
(вывод~\ref{der:10:ground}) даёт
\begin{equation}\label{eq:10:psi100}
  \psi_{100}(r)=\frac{1}{\sqrt{\pi a^3}}\,\ee^{-r/a},\qquad
  a=\frac{\hbar^2}{\mu Ze^2},\qquad E_1=-\frac{\mu Z^2e^4}{2\hbar^2}.
\end{equation}
При $Z=1$, $\mu=\mel$: $a=\aB=0{,}529$~Å, $E_1=-13{,}6$~эВ — те же значения, что в модели Бора.
\emph{Радиальная плотность вероятности} (вероятность найти электрон в сферическом слое $(r,r+\dd r)$):
\begin{equation}\label{eq:10:wr}
  w(r)\,\dd r=|\psi_{100}|^2\,4\pi r^2\dd r=\frac{4r^2}{a^3}\ee^{-2r/a}\dd r ,
\end{equation}
максимальна при $r=a$ — наиболее вероятное расстояние совпадает с боровским радиусом. При
этом $\avg r=\frac32a$, а сама плотность $|\psi|^2$ максимальна в центре ($r=0$). Средние энергии:
$\avg U=-Ze^2/a=2E_1$, $\avg T=-E_1$ (вывод~\ref{der:10:means}) — теорема вириала для кулоновского поля.

\subsection{Классификация состояний (без вывода)}

\begin{principle}{Спектр атома водорода (без вывода)}
Решение радиального уравнения~\eqref{eq:10:radial}, конечное при $r=0$ и убывающее при $r\to\infty$,
существует при
\begin{equation}\label{eq:10:En}
  E_n=-\frac{\mu Z^2e^4}{2\hbar^2n^2}=-\frac{Z^2\Ry}{n^2}\frac{\mu}{\mel},\qquad n=n_r+l+1=1,2,\dots,
\end{equation}
где $n_r=0,1,\dots$ — число узлов $R_{nl}$. При данном $n$: $l=0,1,\dots,n-1$; при данном $l$: $m=-l,\dots,l$.
\end{principle}

\paragraph{Вырождение.} Энергия зависит только от $n$. Число состояний с данным $n$:
\begin{equation}\label{eq:10:deg}
  g_n=\sum_{l=0}^{n-1}(2l+1)=n^2\qquad(2n^2\ \text{с учётом спина}).
\end{equation}
Вырождение по $m$ — следствие сферической симметрии (общее для любого центрального поля).
Вырождение по $l$ — особенность кулоновского поля $1/r$ («случайное» вырождение, связанное с
дополнительной сохраняющейся величиной — вектором Рунге--Ленца). В некулоновском поле, например
для валентного электрона щелочного атома, уровни с разными $l$ расщепляются (квантовый дефект,
\tref{t:03}).

\paragraph{Спектроскопические обозначения.} Состояние записывают как $n$ и букву для $l$:
\begin{center}
\begin{tabular}{@{}lccccc@{}}
\toprule
$n$ & $l=0$ & $l=1$ & $l=2$ & $l=3$ & $g_n$ (без спина)\\
\midrule
1 & $1s$ & & & & 1\\
2 & $2s$ & $2p$ & & & 4\\
3 & $3s$ & $3p$ & $3d$ & & 9\\
4 & $4s$ & $4p$ & $4d$ & $4f$ & 16\\
\bottomrule
\end{tabular}
\end{center}

\paragraph{Радиальные функции нижних состояний} ($Z=1$, $\rho=r/a$, нормировка $\int_0^\infty R^2r^2\dd r=1$):
\[
  R_{10}=\frac{2}{a^{3/2}}\ee^{-\rho},\qquad
  R_{20}=\frac{1}{\sqrt2\,a^{3/2}}\Bigl(1-\frac\rho2\Bigr)\ee^{-\rho/2},\qquad
  R_{21}=\frac{1}{2\sqrt6\,a^{3/2}}\,\rho\,\ee^{-\rho/2}.
\]
Общие свойства: $R_{nl}\propto r^l$ при $r\to0$ (центробежный барьер вытесняет электрон из центра,
поэтому лишь $s$-электроны имеют ненулевую плотность в ядре); $R_{nl}\propto\ee^{-Zr/na}$ при $r\to\infty$;
характерный размер растёт как $n^2a/Z$ (рис.~\ref{fig:10:radial}).

\begin{figure}[htbp]
\centering
\begin{tikzpicture}
\begin{axis}[width=0.75\linewidth, height=5.2cm, xlabel={$r/\aB$}, ylabel={$r^2R_{nl}^2$ (в ед. $1/\aB$)},
   xmin=0, xmax=16, ymin=0, ymax=0.58, samples=300,
   legend style={font=\small, draw=none}, legend pos=north east]
  \addplot[very thick, accent, domain=0:16] {4*x^2*exp(-2*x)};
  \addlegendentry{$1s$}
  \addplot[very thick, keycol, domain=0:16] {0.5*x^2*(1-x/2)^2*exp(-x)};
  \addlegendentry{$2s$}
  \addplot[very thick, dercol, dashed, domain=0:16] {(1/24)*x^4*exp(-x)};
  \addlegendentry{$2p$}
\end{axis}
\end{tikzpicture}
\caption{Радиальные плотности вероятности $w(r)=r^2R_{nl}^2$ для состояний $1s$, $2s$, $2p$ атома водорода
(расчёт по приведённым формулам). Максимум $1s$ — при $r=\aB$, максимум $2p$ — при $r=4\aB$ (как радиус
второй боровской орбиты); у $2s$ — один радиальный узел при $r=2\aB$.}
\label{fig:10:radial}
\end{figure}

\paragraph{Правила отбора и спектр.} Для электрических дипольных переходов $\Delta l=\pm1$,
$\Delta m=0,\pm1$; $\Delta n$ — любое. Частоты переходов даются той же формулой Ридберга, что и в
модели Бора. Состояние $2s$ не может распасться в $1s$ с испусканием одного фотона
($\Delta l=0$) — оно метастабильно (время жизни $\approx0{,}12$~с, распад с испусканием двух фотонов).

\section{Математические выводы}

\begin{derivation}[label=der:10:ground]{энергия и волновая функция основного состояния}
\emph{Шаг 1.} Для $l=0$ функция зависит только от $r$, и по~\eqref{eq:10:lap}
$\nabla^2\psi=\psi''+\frac2r\psi'$. Уравнение~\eqref{eq:10:sse}:
\[
  -\frac{\hbar^2}{2\mu}\Bigl(\psi''+\frac2r\psi'\Bigr)-\frac{Ze^2}{r}\psi=E\psi .
\]
\emph{Шаг 2 (асимптотика).} При $r\to\infty$ остаётся $\psi''\approx-\frac{2\mu E}{\hbar^2}\psi$; при $E<0$
убывающее решение $\psi\sim\ee^{-\varkappa r}$, $\varkappa=\sqrt{-2\mu E}/\hbar$. Пробуем $\psi=C\ee^{-r/a}$.

\emph{Шаг 3 (подстановка).} $\psi'=-\psi/a$, $\psi''=\psi/a^2$:
\[
  -\frac{\hbar^2}{2\mu}\Bigl(\frac{1}{a^2}-\frac{2}{ar}\Bigr)-\frac{Ze^2}{r}=E
  \quad\Longleftrightarrow\quad
  \Bigl(\frac{\hbar^2}{\mu a}-Ze^2\Bigr)\frac1r=E+\frac{\hbar^2}{2\mu a^2}.
\]
\emph{Шаг 4.} Равенство при всех $r$ возможно, только если обе части равны нулю:
\[
  \frac{\hbar^2}{\mu a}=Ze^2\ \Rightarrow\ a=\frac{\hbar^2}{\mu Ze^2};\qquad
  E=-\frac{\hbar^2}{2\mu a^2}=-\frac{\hbar^2}{2\mu}\cdot\frac{\mu^2Z^2e^4}{\hbar^4}=-\frac{\mu Z^2e^4}{2\hbar^2}.
\]
\emph{Шаг 5 (нормировка).} С интегралом $\int_0^\infty r^n\ee^{-\beta r}\dd r=n!/\beta^{n+1}$:
\[
  \int|\psi|^2\dd V=|C|^2\int_0^\infty\ee^{-2r/a}4\pi r^2\dd r=4\pi|C|^2\frac{2!}{(2/a)^3}=\pi a^3|C|^2=1
  \ \Rightarrow\ C=\frac{1}{\sqrt{\pi a^3}}.
\]
Функция не имеет узлов, поэтому это действительно основное состояние.
\end{derivation}

\begin{derivation}[label=der:10:means]{наиболее вероятное расстояние, $\avg r$, средние энергии}
\emph{Наиболее вероятное $r$.} $w(r)\propto r^2\ee^{-2r/a}$; $\dfrac{\dd w}{\dd r}\propto\bigl(2r-\dfrac{2r^2}{a}\bigr)\ee^{-2r/a}=0$ $\Rightarrow$ $r=a$.

\emph{Среднее расстояние.}
$\avg r=\dfrac{4}{a^3}\displaystyle\int_0^\infty r^3\ee^{-2r/a}\dd r=\dfrac{4}{a^3}\cdot\dfrac{3!}{(2/a)^4}=\dfrac{4\cdot6\,a^4}{16\,a^3}=\dfrac32a .$

\emph{Средняя потенциальная энергия.}
$\avg{1/r}=\dfrac{4}{a^3}\displaystyle\int_0^\infty r\,\ee^{-2r/a}\dd r=\dfrac{4}{a^3}\cdot\dfrac{a^2}{4}=\dfrac1a$, поэтому
$\avg U=-Ze^2\avg{1/r}=-\dfrac{Ze^2}{a}=-\dfrac{\mu Z^2e^4}{\hbar^2}=2E_1$.

\emph{Средняя кинетическая энергия.} $\avg T=E_1-\avg U=-E_1=\dfrac{\mu Z^2e^4}{2\hbar^2}>0$. Соотношение
$2\avg T=-\avg U$ — квантовая теорема вириала для потенциала $\propto1/r$.

\emph{Вероятность найти электрон дальше $r=2a$} (оценка «размытости» облака):
$\int_{2a}^\infty w\,\dd r=\ee^{-4}(1+4+8)=13\ee^{-4}\approx0{,}24$.
\end{derivation}

\begin{derivation}[label=der:10:deg]{кратность вырождения $n^2$}
$g_n=\sum_{l=0}^{n-1}(2l+1)=2\cdot\frac{(n-1)n}{2}+n=n^2-n+n=n^2$. С учётом двух проекций спина — $2n^2$:
это ёмкость электронной оболочки с данным $n$ (2, 8, 18, 32), важная для периодической системы (\tref{t:13}).
\end{derivation}

\begin{derivation}[label=der:10:2p]{проверка: $R_{21}\propto r\ee^{-r/2a}$ даёт $E_2=E_1/4$}
Для $l=1$ в уравнении~\eqref{eq:10:radial} ($Z=1$) положим $R=r\ee^{-r/b}$. Тогда
$R'=(1-r/b)\ee^{-r/b}$, $\frac{1}{r^2}(r^2R')'=\bigl(\frac2r-\frac4b+\frac{r}{b^2}\bigr)\ee^{-r/b}$. Подстановка и деление на $\ee^{-r/b}$:
\[
  -\frac{\hbar^2}{2\mu}\Bigl(\frac2r-\frac4b+\frac{r}{b^2}\Bigr)+\frac{\hbar^2}{\mu r}-e^2=Er .
\]
Члены $\propto1/r$ сокращаются автоматически (в этом роль центробежного слагаемого). Члены без $r$:
$\frac{2\hbar^2}{\mu b}=e^2$ $\Rightarrow$ $b=2\aB$. Члены $\propto r$: $E=-\frac{\hbar^2}{2\mu b^2}=-\frac{\hbar^2}{8\mu\aB^2}=\frac{E_1}{4}$ — в согласии с $E_n=E_1/n^2$
при $n=2$. Максимум $w=r^2R^2\propto r^4\ee^{-r/\aB}$ — при $r=4\aB$.
\end{derivation}

\section{Итоговые формулы}

\keyformula{\op H=-\frac{\hbar^2}{2\mu}\nabla^2-\frac{Ze^2}{r},\qquad \psi_{nlm}=R_{nl}(r)Y_{lm}(\theta,\varphi)}
 {$\mu$ — приведённая масса, $Ze$ — заряд ядра, $Y_{lm}$ — сферические функции.}
 {нерелятивистское приближение, точечное ядро, без спина.}
 {постановка задачи; разделение на радиальную и угловую части.}

\keyformula{\psi_{100}=\frac{\ee^{-r/a}}{\sqrt{\pi a^3}},\quad a=\frac{\hbar^2}{\mu Ze^2},\quad E_1=-\frac{\mu Z^2e^4}{2\hbar^2},\quad r_{\text{вер}}=a,\ \avg r=\tfrac32a}
 {$a$ — боровский радиус (для $Z=1$, $\mu=\mel$: $0{,}529$~Å), $E_1=-13{,}6$~эВ.}
 {основное состояние ($n=1$, $l=m=0$).}
 {распределение электрона, размер атома, энергия ионизации.}

\keyformula{E_n=-\frac{Z^2\Ry}{n^2},\quad n=n_r+l+1,\quad l\le n-1,\quad |m|\le l,\quad g_n=n^2\ (2n^2)}
 {$n$, $l$, $m$ — главное, орбитальное, магнитное квантовые числа; $n_r$ — число радиальных узлов.}
 {кулоновское поле (без вывода).}
 {классификация состояний, кратность вырождения, ёмкость оболочек.}

\section{Физическая интерпретация}

\textbf{Сравнение с моделью Бора.} Уровни $E_n$ совпадают, но физическая картина иная: вместо
орбит — распределения вероятности; основное состояние имеет $L=0$ (у Бора $L=\hbar$); боровский
радиус — наиболее вероятное, а не точное расстояние. Совпадение радиусов $n^2\aB$ с максимумами
$w(r)$ имеет место только для состояний с $l=n-1$ (без радиальных узлов).

\textbf{Предельные случаи и масштабы.} Размер $\propto n^2/Z$, энергия связи $\propto Z^2/n^2$. Электроны
$s$-состояний имеют конечную плотность в ядре $|\psi(0)|^2=Z^3/(\pi n^3a^3)$ — это определяет
сверхтонкое расщепление, вероятность электронного захвата (\tref{t:20}) и изотопический «объёмный» сдвиг.

\textbf{Эксперимент.} Спектр водорода — один из наиболее точно измеренных: формула~\eqref{eq:10:En}
выполняется до $\sim10^{-5}$; оставшиеся расхождения (тонкая структура, лэмбовский сдвиг) объясняются
релятивистской теорией и квантовой электродинамикой.

\section{Что необходимо знать на экзамене}

\begin{examcheck}
\checkhead{Определения}
\begin{itemize}
 \item гамильтониан атома водорода, приведённая масса; квантовые числа $n,l,m,m_s$ и их допустимые значения;
 \item эффективный потенциал и центробежный барьер; радиальная плотность вероятности.
\end{itemize}
\checkhead{Формулы}
\begin{itemize}
 \item $\psi_{100}=\ee^{-r/a}/\sqrt{\pi a^3}$, $a=\hbar^2/\mu Ze^2$, $E_1=-\mu Z^2e^4/2\hbar^2$;
 \item $E_n=-Z^2\Ry/n^2$, $n=n_r+l+1$, вырождение $n^2$ ($2n^2$) — без вывода; обозначения $1s$, $2p$, \ldots
\end{itemize}
\checkhead{Выводы}
\begin{itemize}
 \item подстановка $C\ee^{-r/a}$ в уравнение Шредингера (шаги 1--4), нормировка;
 \item $r_{\text{вер}}=a$, $\avg r=\frac32a$, $\avg U=2E$, $\avg T=-E$; сумма $\sum(2l+1)=n^2$.
\end{itemize}
\checkhead{Интерпретации}
\begin{itemize}
 \item происхождение вырождения по $m$ и по $l$; почему у щелочных атомов вырождение по $l$ снято;
 \item отличия от модели Бора; плотность $s$-электронов в ядре; правила отбора и метастабильность $2s$.
\end{itemize}
\end{examcheck}
